% ============================================================================
% CAPÍTULO 6: DISCUSIÓN
% ============================================================================

\chapter{Discusión}
\label{cap:discusion}

% ----------------------------------------------------------------------------
% INTRODUCCIÓN
% ----------------------------------------------------------------------------
En este capítulo se interpretan los resultados presentados en el \Cref{cap:resultados}, analizando su significado, implicaciones y limitaciones. Se contrastan los hallazgos con la literatura revisada en el \Cref{cap:estado_arte}.

% ----------------------------------------------------------------------------
% SECCIÓN 6.1: Interpretación de Resultados Principales
% ----------------------------------------------------------------------------
\section{Interpretación de Resultados Principales}
\label{sec:interpretacion}

Los resultados obtenidos (F1-Score = 0.88) superan significativamente a los métodos baseline, demostrando la efectividad del enfoque propuesto.

\subsection{Comparación con la Literatura}

En contraste con trabajos previos \cite{smith2020,jones2021}, nuestro modelo alcanza un desempeño superior...

La \Cref{tab:comparacion_literatura_discusion} sitúa estos resultados en contexto:

\begin{table}[htbp]
    \centering
    \caption{Comparación con trabajos relacionados}
    \label{tab:comparacion_literatura_discusion}
    \begin{tabular}{@{}llcc@{}}
        \toprule
        \textbf{Estudio} & \textbf{Método} & \textbf{F1-Score} & \textbf{Dataset} \\
        \midrule
        Smith et al. (2020)  & Random Forest & 0.81 & 5,000 instancias \\
        Jones \& Lee (2021)  & SVM          & 0.84 & 8,000 instancias \\
        García et al. (2022) & CNN          & 0.86 & 12,000 instancias \\
        \textbf{Este trabajo} & \textbf{Modelo Propuesto} & \textbf{0.88} & \textbf{10,000 instancias} \\
        \bottomrule
    \end{tabular}
\end{table}

\subsection{Factores que Contribuyen al Desempeño}

El desempeño superior puede atribuirse a:

\begin{enumerate}
    \item \textbf{Arquitectura optimizada:} La combinación de capas con dropout reduce el sobreajuste
    \item \textbf{Preprocesamiento robusto:} El manejo cuidadoso de datos faltantes mejora la calidad...
    \item \textbf{Ingeniería de características:} La selección de variables mediante \textit{[técnica]} captura...
\end{enumerate}

% ----------------------------------------------------------------------------
% SECCIÓN 6.2: Análisis de Casos Específicos
% ----------------------------------------------------------------------------
\section{Análisis de Casos Específicos}
\label{sec:casos}

\subsection{Casos de Éxito}

El modelo clasifica correctamente casos como...

\subsection{Casos de Error}

Los errores más comunes ocurren cuando...

Por ejemplo, en la \Cref{fig:caso_error}, se observa que...

\begin{figure}[htbp]
    \centering
    \includegraphics[width=0.7\textwidth]{figuras/caso_error_analisis.png}
    \caption{Ejemplo de caso mal clasificado. La instancia tiene características ambiguas...}
    \label{fig:caso_error}
\end{figure}

Estos errores sugieren que el modelo podría mejorar con...

% ----------------------------------------------------------------------------
% SECCIÓN 6.3: Implicaciones Prácticas
% ----------------------------------------------------------------------------
\section{Implicaciones Prácticas}
\label{sec:implicaciones}

Los resultados tienen implicaciones importantes para:

\begin{itemize}
    \item \textbf{Área de aplicación 1:} El modelo puede ser utilizado para...
    \item \textbf{Área de aplicación 2:} Los hallazgos sugieren que...
    \item \textbf{Área de aplicación 3:} La metodología propuesta facilita...
\end{itemize}

% ----------------------------------------------------------------------------
% SECCIÓN 6.4: Limitaciones del Estudio
% ----------------------------------------------------------------------------
\section{Limitaciones del Estudio}
\label{sec:limitaciones}

Es importante reconocer las siguientes limitaciones:

\begin{enumerate}
    \item \textbf{Tamaño del dataset:} Con 10,000 instancias, el modelo podría mejorar con más datos...
    
    \item \textbf{Generalización:} Los resultados se obtuvieron con datos de \textit{[contexto específico]}, por lo que la generalización a otros contextos requiere validación adicional.
    
    \item \textbf{Recursos computacionales:} El entrenamiento requiere \textit{[tiempo/recursos]}, lo cual puede limitar su aplicación en...
    
    \item \textbf{Interpretabilidad:} Como modelo de red neuronal, la interpretación de decisiones individuales es limitada...
\end{enumerate}

% ----------------------------------------------------------------------------
% SECCIÓN 6.5: Validez y Confiabilidad
% ----------------------------------------------------------------------------
\section{Validez y Confiabilidad}
\label{sec:validez}

\subsection{Validez Interna}

El uso de validación cruzada y la evaluación con múltiples métricas aseguran...

\subsection{Validez Externa}

Los resultados son generalizables a \textit{[contextos similares]} debido a...

\subsection{Confiabilidad}

La reproducibilidad del estudio está garantizada por:
\begin{itemize}
    \item Descripción detallada de la metodología
    \item Especificación de hiperparámetros
    \item Disponibilidad de código fuente\footnote{Disponible en: https://github.com/usuario/proyecto}
\end{itemize}

% ----------------------------------------------------------------------------
% SECCIÓN 6.6: Relación con los Objetivos
% ----------------------------------------------------------------------------
\section{Cumplimiento de Objetivos}
\label{sec:cumplimiento_objetivos}

Los resultados demuestran que se han alcanzado los objetivos planteados:

\begin{itemize}
    \item \textbf{Objetivo General:} Se desarrolló exitosamente un sistema que...
    \item \textbf{OE1:} El análisis de la literatura (Capítulo 3) identificó...
    \item \textbf{OE2:} El diseño del modelo (Capítulo 4) incorporó...
    \item \textbf{OE3:} La implementación (Capítulo 4) demostró ser...
    \item \textbf{OE4:} La evaluación (Capítulo 5) confirmó que...
\end{itemize}

% ============================================================================
% CONSEJOS:
% ============================================================================
% - INTERPRETA los resultados (no solo los repitas)
% - Relaciona con la literatura: ¿confirman o contradicen trabajos previos?
% - Explica el "por qué" de los resultados
% - Sé honesto con las limitaciones
% - Conecta hallazgos con implicaciones prácticas
% - Sugiere mejoras futuras
% ============================================================================


